\ifdefined\gkver\else\def\gkver{student}\fi
\documentclass[\gkver]{gaokaozhenti}
\begin{document}

\chapter{2023年6月浙江}

\begin{enumerate}
\item
%% number: 1
%% paperName: 2023年6月浙江高考第 1 题 3 分
%% typeId: 1
%% type: 单选题
%% chapter: 功和能
%% point: 功率
%% method: 无
%% score: 3
%% degree: 950
%% duplicateId: 0
%% body:
下列四组物理量中均为标量的是\xzanswer{B}
\fourchoices[answer=B]
{电势  电场强度}
{热量  功率}
{动量  动能}
{速度  加速度}
%% answer: B
%% memo:
\memoanswer{
A．电势只有大小没有方向，是标量；电场强度既有大小又有方向，是矢量，故 A 错误；

B．热量和功率都是只有大小没有方向，都是标量，故 B 正确；

C．动量既有大小又有方向，是矢量；动能只有大小没有方向，是标量，故 C 错误；

D．速度和加速度都是既有大小又有方向的物理量，是矢量，故 D 错误。

故选 B。
}

\item
%% number: 2
%% paperName: 2023年6月浙江高考第 2 题 3 分
%% typeId: 1
%% type: 单选题
%% chapter: 运动和力
%% point: 牛顿第一定律
%% method: 无
%% score: 3
%% degree: 950
%% duplicateId: 0
%% body:
在足球运动中，足球入网如图所示，则\xzanswer{B}
\onepicture{figs/02.jpg}
\fourchoices[answer=B]
{踢香蕉球时足球可视为质点}
{足球在飞行和触网时惯性不变}
{足球在飞行时受到脚的作用力和重力}
{触网时足球对网的力大于网对足球的力}
%% answer: B
%% memo:
\memoanswer{
A．在研究如何踢出“香蕉球”时，需要考虑踢在足球上的位置与角度，所以不可以把足球看作质点，故 A 错误；

B．惯性只与质量有关，足球在飞行和触网时质量不变，则惯性不变，故 B 正确；

C．足球在飞行时脚已经离开足球，故在忽略空气阻力的情况下只受重力，故 C 错误；

D．触网时足球对网的力与网对足球的力是相互作用力，大小相等，故 D 错误。

故选 B。
}

\item
%% number: 3
%% paperName: 2023年6月浙江高考第 3 题 3 分
%% typeId: 1
%% type: 单选题
%% chapter: 功和能
%% point: 功能中的图象
%% method: 无
%% score: 3
%% degree: 850
%% duplicateId: 0
%% body:
铅球被水平推出后的运动过程中，不计空气阻力，下列关于铅球在空中运动时的加速度大小 $a$、速度大小 $v$、动能 $E_{k}$ 和机械能 $E $随运动时间 $t$ 的变化关系中，正确的是\xzanswer{D}
\onepicture{\includesvg{figs/03}}
%% answer: D
%% memo:
\memoanswer{
A．由于不计空气阻力，铅球被水平推出后只受重力作用，加速度等于重力加速度，不随时间改变，故 A 错误；

B．铅球被水平推出后做平抛运动，竖直方向有  $v_{y}=gt$，则抛出后速度大小为  $v=\sqrt{v_{0}^{2}+\left(gt\right)^{2}}$，可知速度大小与时间不是一次函数关系，故 B 错误；

C. 铅球抛出后的动能

 $E_{k}=\frac{1}{2}mv^{2}=\frac{1}{2}m\left[v_{0}^{2}+\left(gt\right)^{2}\right]$

可知动能与时间不是一次函数关系，故 C 错误；

D．铅球水平抛出后由于忽略空气阻力，所以抛出后铅球机械能守恒，故 D 正确。故选 D。
}

\item
%% number: 4
%% paperName: 2023年6月浙江高考第 4 题 3 分
%% typeId: 1
%% type: 单选题
%% chapter: 曲线运动
%% point: 曲线运动条件
%% method: 无
%% score: 3
%% degree: 850
%% duplicateId: 0
%% body:
图为“玉兔二号”巡视器在月球上从 O 处行走到 B 处的照片，轨迹 OA 段是直线，AB 段是曲线，巡视器质量为 135 kg，则巡视器\xzanswer{B}
\onepicture{figs/04.jpg}
\fourchoices[answer=B]
{受到月球的引力为 1 350 N}
{在 AB 段运动时一定有加速度}
{OA 段与 AB 段的平均速度方向相同}
{从 O 到 B 的位移大小等于 OAB 轨迹长度}
%% answer: B
%% memo:
\memoanswer{
A．在月球上的 $g$ 与地球不同，故质量为 135 kg 的巡视器受到月球的引力不是1 350 N，故 A 错误；

B．由于在 AB 段运动时做曲线运动，速度方向一定改变，一定有加速度，故 B 正确；

C．平均速度的方向与位移方向相同，由图可知 OA 段与 AB 段位移方向不同，故平均速度方向不相同，故 C 错误；

D．根据位移的定义可知从 O 到 B 的位移大小等于 OB 的连线长度，故 D 错误。

故选 B。
}

\item
%% number: 5
%% paperName: 2023年6月浙江高考第 5 题 3 分
%% typeId: 1
%% type: 单选题
%% chapter: 原子和原子核
%% point: 结合能
%% method: 无
%% score: 3
%% degree: 950
%% duplicateId: 0
%% body:
“玉兔二号”装有核电池，不惧漫长寒冷的月夜。核电池将 $^{238}_{94}$Pu 衰变释放的核能一部分转换成电能。$^{238}_{94}$Pu 的衰变方程为 $^{238}_{94}$Pu $\to ^{X}_{92}$U + $^{4}_{2}$He，则\xzanswer{C}
\fourchoices[answer=C]
{衰变方程中的 X 等于 233}
{$^{4}_{2}$He 的穿透能力比 $\gamma$ 射线强}
{$^{238}_{94}$Pu 比 $^{X}_{92}$U 的比结合能小}
{月夜的寒冷导致 $^{238}_{94}$Pu 的半衰期变大}
%% answer: C
%% memo:
\memoanswer{
A．根据质量数和电荷数守恒可知，衰变方程为

$^{238}_{94}$Pu $\to ^{234}_{92}$U + $^{4}_{2}$He

即衰变方程中的 X = 234，故 A 错误；

B．$^{4}_{2}$He 是 $\alpha$ 粒子，穿透能力比 $\gamma$ 射线弱，故 B 错误；

C．比结合能越大越稳定，由于 $^{238}_{94}$Pu 衰变成为了 $^{234}_{92}$U，故 $^{234}_{92}$U 比 $^{238}_{94}$Pu 稳定，即 $^{238}_{94}$Pu 比 $^{234}_{92}$U 的比结合能小，故 C 正确；

D．半衰期由原子核本身决定的，与温度等外部因素无关，故 D 错误。

故选 C。
}

\item
%% number: 6
%% paperName: 2023年6月浙江高考第 6 题 3 分
%% typeId: 1
%% type: 单选题
%% chapter: 力物体的平衡
%% point: 三力平衡
%% method: 无
%% score: 3
%% degree: 850
%% duplicateId: 0
%% body:
如图所示，水平面上固定两排平行的半圆柱体，重为 $G$ 的光滑圆柱体静置其上，a、b 为相切点，$\angle$aOb = 90${^\circ}$，半径 Ob 与重力的夹角为 37${^\circ}$。已知 sin37${^\circ}$ = 0.6，cos37${^\circ}$ = 0.8，则圆柱体受到的支持力 $F_{a}$、$F_{b}$ 大小为\xzanswer{D}
\onepicture{figs/06.png}
\fourchoices[answer=D]
{$F_{a}$ = 0.6$G$，$F_{b}$ = 0.4$G$}
{$F_{a}$ = 0.4$G$，$F_{b}$ = 0.6$G$}
{$F_{a}$ = 0.8$G$，$F_{b}$ = 0.6$G$}
{$F_{a}$ = 0.6$G$，$F_{b}$ = 0.8$G$}
%% answer: D
%% memo:
\memoanswer{
\onepicture[w=6cm]{figs/06_an.png}

对光滑圆柱体受力分析如图。由题意有

$F_{a} = G$sin37${^\circ}$ = 0.6$G$，$F_{b} = G$cos37${^\circ}$ = 0.8$G$

故选 D。
}

\item
%% number: 7
%% paperName: 2023年6月浙江高考第 7 题 3 分
%% typeId: 1
%% type: 单选题
%% chapter: 电磁感应
%% point: 远距离输电
%% method: 无
%% score: 3
%% degree: 750
%% duplicateId: 0
%% body:
我国 1 100 kV 特高压直流输电工程的送电端用“整流”设备将交流变换成直流，用户端用“逆变”设备再将直流变换成交流。下列说法正确的是\xzanswer{A}
\fourchoices[answer=A]
{送电端先升压再整流}
{用户端先降压再变交流}
{1 100 kV 是指交流电的最大值}
{输电功率由送电端电压决定}
%% answer: A
%% memo:
\memoanswer{
AB．升压和降压都需要在交流的时候才能进行，故送电端应该先升压再整流，用户端应该先变交流再降压，故 A 正确，B 错误；

C．1 100 kV 指的是交流电的有效值，故 C 错误；

D．输电的功率是由用户端负载的总功率来决定的，故 D 错误。

故选 A。
}

\item
%% number: 8
%% paperName: 2023年6月浙江高考第 8 题 3 分
%% typeId: 1
%% type: 单选题
%% chapter: 电场
%% point: 带电粒子的圆周运动
%% method: 无
%% score: 3
%% degree: 750
%% duplicateId: 0
%% body:
某带电粒子转向器的横截面如图所示，转向器中有辐向电场。粒子从 M 点射入，沿着由半径分别为 $R_{1}$ 和 $R_{2}$ 的圆弧平滑连接成的虚线（等势线）运动，并从虚线上的 N 点射出，虚线处电场强度大小分别为 $E_{1}$ 和 $E_{2}$，则 $R_{1}$、$R_{2}$ 和 $E_{1}$、$E_{2}$ 应满足\xzanswer{A}
\onepicture{figs/08.png}
\fourchoices[answer=A]
{$\frac{{{E_1}}}{{{E_2}}} = \frac{{{R_2}}}{{{R_1}}}$}
{$\frac{{{E_1}}}{{{E_2}}} = \frac{{R_1^2}}{{R_2^2}}$}
{$\frac{{{E_1}}}{{{E_2}}} = \frac{{{R_1}}}{{{R_2}}}$}
{$\frac{{{E_1}}}{{{E_2}}} = \frac{{R_2^2}}{{R_1^2}}$}
%% answer: A
%% memo:
\memoanswer{
带电粒子在电场中做匀速圆周运动，电场力提供向心力，则有

 $qE_{1}=m\frac{v^{2}}{R_{1}}\quad qE_{2}=m\frac{v^{2}}{R_{2}}$

联立可得  $\frac{E_{1}}{E_{2}} = \frac{R_{2}}{R_{1}}$

故选 A。
}

\item
%% number: 9
%% paperName: 2023年6月浙江高考第 9 题 3 分
%% typeId: 1
%% type: 单选题
%% chapter: 曲线运动
%% point: 天体运动
%% method: 无
%% score: 3
%% degree: 450
%% duplicateId: 0
%% body:
木星的卫星中，木卫一、木卫二、木卫三做圆周运动的周期之比为 1$:$2$:$4。木卫三周期为 $T$，公转轨道半径是月球绕地球轨道半径 $r$ 的 $n$ 倍。月球绕地球公转周期为 $T_{0}$，则\xzanswer{D}
\fourchoices[answer=D]
{木卫一轨道半径为 $\frac{n}{{16}} r$}
{木卫二轨道半径为 $\frac{n}{{2}} r$}
{周期 $T$ 与 $T_{0}$ 之比为 $\sqrt {{n^3}} $}
{木星质量与地球质量之比为 $\frac{{T_0^2}}{{{T^2}}}n^{3}$}
%% answer: D
%% memo:
\memoanswer{
根据题意可得，木卫 3 的轨道半径为  $r_{3}=nr$

AB. 根据万有引力提供向心力  $G\frac{Mm}{R^{2}}=m\frac{4\pi^{2}}{T^{2}}R$

可得  $R=\sqrt[3]{\frac{GMT^{2}}{4\pi^{2}}}$

木卫一、木卫二、木卫三做圆周运动的周期之比为 1:2:4，可得

木卫一轨道半径为  $r_{1}=\frac{nr}{\sqrt[3]{16}}$，木卫二轨道半径为  $r_{2}=\frac{nr}{\sqrt[3]{4}}$

故 AB 错误；

C. 木卫三围绕的中心天体是木星，月球的围绕的中心天体是地球，根据题意无法求出周期  $T$ 与  $T_{0}$ 之比，故 C 错误；

D. 根据万有引力提供向心力，分别有

 $G\frac{M_{\text{木}}m}{(nr)^{2}}=m\frac{4\pi^{2}}{T^{2}}nr$， $G\frac{M_{\text{地}}m}{r^{2}}=m\frac{4\pi^{2}}{T_{0}^{2}}r$

联立可得  $\frac{M_{\text{木}}}{M_{\text{地}}} = \frac{T_{0}^{2}}{T^{2}} n^{3}$

故 D 正确。

故选 D。
}

\item
%% number: 10
%% paperName: 2023年6月浙江高考第 10 题 3 分
%% typeId: 1
%% type: 单选题
%% chapter: 电磁感应
%% point: 电磁感应能量分析
%% method: 无
%% score: 3
%% degree: 550
%% duplicateId: 0
%% body:
如图所示，质量为 $M$、电阻为 $R$、长为 $L$ 的导体棒，通过两根长均为 $l$、质量不计的导电细杆连在等高的两固定点上，固定点间距也为 $L$。细杆通过开关 S 可与直流电源 $E_{0}$ 或理想二极管串接。在导体棒所在空间存在磁感应强度方向竖直向上、大小为 $B$ 的匀强磁场，不计空气阻力和其它电阻。开关 S 接 1，当导体棒静止时，细杆与竖直方向的夹角固定点 $\theta =  \frac{\pi }{4}$；然后开关 S 接 2，棒从右侧开始运动完成一次振动的过程中\xzanswer{C}
\onepicture{figs/10.png}
\fourchoices[answer=C]
{电源电动势 $E_{0} = \frac{{\sqrt 2 MgR}}{{2BL}}$}
{棒消耗的焦耳热 $Q$ = (1 $- \frac{{\sqrt 2 }}{2}$)$Mgl$}
{从左向右运动时，最大摆角小于 $\frac{\pi }{4}$}
{棒两次过最低点时感应电动势大小相等}
%% answer: C
%% memo:
\memoanswer{
\onepicture[w=6cm]{\includesvg{figs/10_an}}

A．当开关接1时，对导体棒受力分析如图所示。

根据几何关系可得 $Mg = BIL$，解得 $I = \frac{{Mg}}{{BL}}$

根据欧姆定律 $I = \frac{{{E_0}}}{R}$ 解得 $E_{0} = \frac{{MgR}}{{2BL}}$

故 A 错误；

根据左手定则可知导体棒从右向左运动时，产生的感应电动势与二极管相反，没有机械能损失，导体棒从左向右运动时，产生的感应电动势与二极管正方向相同，部分机械能转化为焦耳热：

B．当导体棒运动到最低点速度为零时，导体棒损失的机械能转化为焦耳热为 $Q'$ = (1 $- \frac{{\sqrt 2 }}{2}$)$Mgl$，根据楞次定律可知导体棒完成一次振动速度为零时，导体棒高度高于最低点，所以棒消耗的焦耳热 $Q < Q'$ = (1 $- \frac{{\sqrt 2 }}{2}$)$Mgl$。

故 B 错误；

C．根据 B 选项分析可知，导体棒运动过程中，机械能转化为焦耳热，所以从左向右运动时，最大摆角小于 $\frac{\pi }{4}$，故 C 正确；

D．根据 B 选项分析，导体棒第二次经过最低点时的速度小于第一次经过最低点时的速度，根据 $E = BLv$ 可知棒两次过最低点时感应电动势大小不相等，故 D 错误。

故选 C。
}

\item
%% number: 11
%% paperName: 2023年6月浙江高考第 11 题 3 分
%% typeId: 1
%% type: 单选题
%% chapter: 机械波
%% point: 干涉衍射
%% method: 无
%% score: 3
%% degree: 450
%% duplicateId: 0
%% body:
如图所示，置于管口 T 前的声源发出一列单一频率声波，分成两列强度不同的声波分别沿 A、B 两管传播到出口 O。先调节 A、B 两管等长，O 处探测到声波强度为 400 个单位，然后将 A 管拉长 $d$ = 15 cm，在 O 处第一次探测到声波强度最小，其强度为 100 个单位。已知声波强度与声波振幅平方成正比，不计声波在管道中传播的能量损失，则\xzanswer{C}
\onepicture{figs/11.png}
\fourchoices[answer=C]
{声波的波长 $\lambda$ = 15 cm}
{声波的波长 $\lambda$ = 30 cm}
{两声波的振幅之比为 3$:$1}
{两声波的振幅之比为 2$:$1}
%% answer: C
%% memo:
\memoanswer{
CD．分析可知 A、B 两管等长时，声波的振动加强，将 A 管拉长  $d = 15 \, \text{cm}$ 后，两声波在 O 点减弱，根据题意设声波加强时振幅为 20，声波减弱时振幅为 10，则  $A_1 + A_2 = 20$， $A_1 - A_2 = 10$。可得两声波的振幅之比， $A_1 : A_2 = 3 : 1$。故 C 正确，D 错误；

AB．根据振动减弱的条件可得  $\frac{\lambda}{2} = 2d$，解得  $\lambda = 60 \, \text{cm}$。故 AB 错误。

故选 C。
}

\item
%% number: 12
%% paperName: 2023年6月浙江高考第 12 题 3 分
%% typeId: 1
%% type: 单选题
%% chapter: 电场
%% point: 带电粒子的直线运动
%% method: 无
%% score: 3
%% degree: 450
%% duplicateId: 0
%% body:
AB、CD 两块正对的平行金属板与水平面成 30${^\circ}$ 角固定，竖直截面如图所示。两板间距 10 cm，电荷量为 1.0$\times$10$^{-8}$ C、质量为 3.0$\times$10$^{-4}$ kg 的小球用长为 5 cm 的绝缘细线悬挂于 A 点。闭合开关 S，小球静止时，细线与 AB 板夹角为 30${^\circ}$；剪断细线，小球运动到 CD 板上的 M 点（未标出），则\xzanswer{B}
\onepicture{figs/12.png}
\fourchoices[answer=B]
{MC 距离为 5$\sqrt 3 $ cm}
{电势能增加了 $\frac{3}{4}\sqrt 3 \times$10$^{-4}$ J}
{电场强度大小为 $\sqrt 3 \times$10$^{4}$ $\UNC$}
{减小 $R$ 的阻值，MC 的距离将变大}
%% answer: B
%% memo:
\memoanswer{
\onepicture[w=6cm]{\includesvg{figs/12_an}}

A．根据平衡条件和几何关系，对小球受力分析如图所示

根据几何关系可得

$T = qE$

$T$sin60${^\circ} + qE$sin60${^\circ} = mg$

联立解得

$T = qE = \times$10$^{-3}$ N

\onepicture[w=6cm]{figs/12_an2.png}

剪断细线，小球做匀加速直线运动，如图所示

根据几何关系可得 $L_{MC} = d$tan60${^\circ}$ = 10$\sqrt 3 $ cm

故 A 错误；

B．根据几何关系可得小球沿着电场力方向的位移 $x$ = (10 - 5sin30${^\circ}$)cm = 7.5 cm

与电场力方向相反，电场力做功为

$W_{\text{电}} = -qEx = - \frac{3}{4}\sqrt 3 \times$10$^{-4}$ J

则小球的电势能增加 $\frac{3}{4}\sqrt 3 \times$10$^{-4}$ J，故 B 正确；

C．电场强度的大小 $E = \frac{{qE}}{q} = \sqrt 3 \times$10$^{5}$ $\UNC$。故 C 错误；

D．减小 $R$ 的阻值，极板间的电势差不变，极板间的电场强度不变，所以小球的运动不会发生改变，MC 的距离不变，故 D 错误。

故选 B。
}

\item
%% number: 13
%% paperName: 2023年6月浙江高考第 13 题 3 分
%% typeId: 1
%% type: 单选题
%% chapter: 光学
%% point: 全反射
%% method: 无
%% score: 3
%% degree: 350
%% duplicateId: 0
%% body:
在水池底部水平放置三条细灯带构成的等腰直角三角形发光体，直角边的长度为 0.9 m，水的折射率 $n = \frac{4}{3}$，细灯带到水面的距离 $h = \frac{{\sqrt 7 }}{{10}}$ m，则有光射出的水面形状（用阴影表示）为\xzanswer{C}
\onepicture{\includesvg{figs/13}}
%% answer: C
%% memo:
\memoanswer{
\onepicture[w=6cm]{figs/13_an.png}



\onepicture[w=6cm]{figs/13_an2.png}
}

\item
%% number: 14
%% paperName: 2023年6月浙江高考第 14 题 3 分
%% typeId: 2
%% type: 多选题
%% chapter: 机械波
%% point: 多普勒效应
%% method: 无
%% score: 3
%% degree: 750
%% duplicateId: 0
%% body:
下列说法正确的是\xzanswer{BD}
\fourchoices[answer=BD]
{热量能自发地从低温物体传到高温物体}
{液体的表面张力方向总是跟液面相切}
{在不同的惯性参考系中，物理规律的形式是不同的}
{当波源与观察者相互接近时，观察者观测到波的频率大于波源振动的频率}
%% answer: BD
%% memo:
\memoanswer{
A．根据热力学第二定律可知热量能不可能自发地从低温物体传到高温物体，故 A 错误；

B．液体的表面张力方向总是跟液面相切，故 B 正确；

C．由狭义相对论的两个基本假设可知，在不同的惯性参考系中，一切物理规律都是相同的，故 C 错误；

D．根据多普勒效应可知当波源与观察者相互接近时，观察者观测到波的频率大于波源振动的频率，故 D 正确。

故选 BD。
}

\item
%% number: 15
%% paperName: 2023年6月浙江高考第 15 题 3 分
%% typeId: 2
%% type: 多选题
%% chapter: 光学
%% point: 光子说
%% method: 无
%% score: 3
%% degree: 450
%% duplicateId: 0
%% body:
有一种新型光电效应量子材料，其逸出功为 $W_{0}$。当紫外光照射该材料时，只产生动能和动量单一的相干光电子束。用该电子束照射间距为 $d$ 的双缝，在与缝相距为 $L$ 的观测屏上形成干涉条纹，测得条纹间距为 $\Delta x$。已知电子质量为 $m$，普朗克常量为 $h$，光速为 $c$，则\xzanswer{AD}
\fourchoices[answer=AD]
{电子的动量 $p_{c} = \frac{{hL}}{{d\Delta x}}$}
{电子的动能 $E_{k} = \frac{{h{L^2}}}{{2m{d^2}\Delta {x^2}}}$}
{光子的能量 $E = W_{0} + \frac{{chL}}{{d\Delta x}}$}
{光子的动量 $p = \frac{{{W_0}}}{c} + \frac{{{h^2}{L^2}}}{{2cm{d^2}\Delta {x^2}}}$}
%% answer: AD
%% memo:
\memoanswer{
根据条纹间距公式  $\Delta x = \frac{L}{d} \lambda$ 可得  $\lambda = \frac{\Delta x d}{L}$

A. 根据  $p_c = \frac{h}{\lambda}$ 可得  $p_c = \frac{hL}{d\Delta x}$。故 A 正确；

B. 根据动能和动量的关系  $E_k = \frac{p^2}{2m}$，结合 A 选项可得  $E_k = \frac{h^2L^2}{2md^2\Delta x^2}$。故 B 错误；

C. 光子的能量  $E = W_0 + E_k = W_0 + \frac{h^2L^2}{2md^2\Delta x^2}$。故 C 错误；

D. 光子的动量  $p = mc$，光子的能量  $E = mc^2$，联立可得  $p = \frac{E}{c}$。则光子的动量  $p = \frac{W_0}{c} + \frac{h^2L^2}{2cmd^2\Delta x^2}$。故 D 正确。

故选 BD。
}

\item
%% number: 16
%% paperName: 2023年6月浙江高考第 16 题 4 分
%% typeId: 4
%% type: 实验题
%% chapter: 曲线运动
%% point: 研究平抛运动实验
%% method: 无
%% score: 4
%% degree: 550
%% duplicateId: 0
%% body:
在“探究平抛运动的特点”实验中
\begin{enumerate}
	\item
	用图 1 装置进行探究，下列说法正确的是\_\_\_\_\_\_\_\_。
	\item
	用图 2 装置进行实验，下列说法正确的是\_\_\_\_\_\_\_\_。
	\item
	A．斜槽轨道 M 必须光滑且其末端水平
	\item
	B．上下调节挡板 N 时必须每次等间距移动
	\item
	C．小钢球从斜槽 M 上同一位置静止滚下
	\item
	用图 3 装置进行实验，竖直挡板上附有复写纸和白纸，可以记下钢球撞击挡板时的点迹。实验时竖直挡板初始位置紧靠斜槽末端，钢球从斜槽上 P 点静止滚下，撞击挡板留下点迹 0，将挡板依次水平向右移动 $x$，重复实验，挡板上留下点迹 1、2、3、4。以点迹 0 为坐标原点，竖直向下建立坐标轴 $y$，各点迹坐标值分别为 $y_{1}$、$y_{2}$、$y_{3}$、$y_{4}$。测得钢球直径为 $d$，则钢球平抛初速度 $v_{0}$ 为\_\_\_\_\_\_\_\_。
	\item
	A．($x + \frac{d}{2}$)$\sqrt {\frac{g}{{2{y_1}}}} $            B．($x + \frac{d}{2}$)$\sqrt {\frac{g}{{{y_2} - {y_1}}}} $
	\item
	C．(3$x - \frac{d}{2}$)$\sqrt {\frac{g}{{2{y_4}}}} $          D．(4$x - \frac{d}{2}$)$\sqrt {\frac{g}{{2{y_4}}}} $
\end{enumerate}
\onepicture[align=right]{figs/16.png}
%% answer: 
\jdanswer{
\begin{enumerate}
	\item
	B

	\item
	C

	\item
	D

\end{enumerate}
}
%% memo:
\memoanswer{
（1）AC．用如图 1 所示的实验装置，只能探究平抛运动竖直分运动的特点，故 AC 错误；

B．在实验过程中，需要改变小锤击打的力度，多次重复实验，故 B 正确。

故选 B。

（2）AC．为了保证小球做平抛运动，需要斜槽末端水平，为了保证小球抛出时速度相等，每一次小球需要静止从同一位置释放，斜槽不需要光滑，故 A 错误，C 正确；

B．上下调节挡板 N 时不必每次等间距移动，故 B 错误。

故选 C。

（3）A．竖直方向，根据 $y_{1} = \frac{1}{2}gt^{2}$

水平方向 $x - \frac{d}{2} = v_{0}t$

联立可得 $v_{0}$ = ($x - \frac{d}{2}$)$\sqrt {\frac{g}{{2{y_1}}}} $

故 A 错误；

B．竖直方向，根据 $\Delta y = y_{2} -$ 2$y_{1} = gt^{2}$

水平方向 $x = v_{0}t$

联立可得 $v_{0} = x \sqrt {\frac{g}{{{y_2} - {2y_1}}}} $

故 B 错误；

CD．竖直方向根据 $y_{4} = \frac{1}{2}gt^{2}$

水平方向 4$x - \frac{d}{2} = v_{0}t$

联立可得 $v_{0}$ = (4$x - \frac{d}{2}$)$\sqrt {\frac{g}{{2{y_4}}}} $

故 D 正确，C 错误。

故选 D。
}

\item
%% number: 16
%% paperName: 2023年6月浙江高考第 16 题 3 分
%% typeId: 4
%% type: 实验题
%% chapter: 力物体的平衡
%% point: 胡克定律
%% method: 无
%% score: 3
%% degree: 750
%% duplicateId: 0
%% body:
如图所示，某同学把 A、B 两根不同的弹簧串接竖直悬挂，探究 A、B 弹簧弹力与伸长量的关系。在 B 弹簧下端依次挂上质量为 $m$ 的钩码，静止时指针所指刻度 $x_{A}$、$x_{B}$ 的数据如表。

\begin{center}
\begin{tabular}{|c|}
\hline
钩码个数 \\ \hline
\end{tabular}
\end{center}

1

2

$\ldots$

$x_{A}$/cm

$\ldots$

$x_{B}$/cm

$\ldots$

钩码个数为 2 时，弹簧 A 的伸长量 $\Delta x_{A}$ = \_\_\_\_\_\_\_\_cm，弹簧 B 的伸长量 $\Delta x_{B}$ = \_\_\_\_\_\_\_\_cm，两根弹簧弹性势能的增加量 $\Delta E_{p}$ = \_\_\_\_\_\_\_\_$mg$($\Delta x_{A} + \Delta x_{B}$)（选填“=”、“<”或“>”）。
\begin{enumerate}
	\item
	75
	\item
	53
	\item
	30
	\item
	45
	\item
	52
	\item
	60
\end{enumerate}
\onepicture[align=right]{figs/16.png}
%% answer: 
\jdanswer{
0.78，1.29，>
}
%% memo:
\memoanswer{
钩码个数为 2 时，弹簧 A 的伸长量

$\Delta x_{A}$ = 8.53 cm $-$ 7.75 cm = 0.78 cm

弹簧 B 的伸长量

$\Delta x_{B}$ = 18.52 cm - 16.45 cm $-$ 0.78 cm = 1.29 cm

根据系统机械能守恒定律可知两根弹簧的重力势能减少量和钩码减少的重力势能等于弹簧增加的弹性势能

$\Delta E_{p} > mg$($\Delta x_{A} + \Delta x_{B}$)
}

\item
%% number: 16
%% paperName: 2023年6月浙江高考第 16 题 5 分
%% typeId: 4
%% type: 实验题
%% chapter: 电路
%% point: 测电动势与内阻
%% method: 无
%% score: 5
%% degree: 650
%% duplicateId: 0
%% body:
在“测量干电池的电动势和内阻”实验中
\begin{enumerate}
	\item
	部分连线如图 1 所示，导线 a 端应连接到\_\_\_\_\_\_\_\_（选填“A”、“B”、“C”或“D”）接线柱上。正确连接后，某次测量中电压表指针位置如图 2 所示，其示数为\_\_\_\_\_\_\_\_V。
	\item
	测得的 7 组数据已标在如图 3 所示 $U - I$ 坐标系上，用作图法求干电池的电动势 $E$ = \_\_\_\_\_\_\_\_V 和内阻 $r$ = \_\_\_\_\_\_\_\_$\Omega$。（计算结果均保留两位小数）
\end{enumerate}
\threepicture[labelA=20236月浙江16a, labelB=20236月浙江16b, labelC=20236月浙江16c, align=right]
{figs/16a.png}
{figs/16b.png}
{figs/16c.png}
%% answer: 
\jdanswer{
\begin{enumerate}
	\item
	B，1.20

	\item
	1.50，1.04

\end{enumerate}
}
%% memo:
\memoanswer{
（1）电压表测量的电压应为滑动变阻器接入电路中电阻丝两端的电阻，开关应能控制电路，所以导线 a 端应连接到 B 处；

干电池电动势约为 1.5 V，电压表选择 0 \textasciitilde{} 3 V 量程，分度值为 0.1 V，题图中电压表读数为 1.20 V；

（2）作出 $U - I$ 如图所示

\onepicture[w=6cm]{figs/16_an.png}

根据闭合电路欧姆定律 $U = E - Ir$

可知 $U - I$ 图像纵轴截距为电源电动势，可得 $E$ = 1.50 V；$U - I$ 图像斜率的绝对值等于电源内阻 $r$ = 1.04 $\Omega$。
}

\item
%% number: 16
%% paperName: 2023年6月浙江高考第 16 题 2 分
%% typeId: 4
%% type: 实验题
%% chapter: 交流电
%% point: 变压器
%% method: 无
%% score: 2
%% degree: 750
%% duplicateId: 0
%% body:
以下实验中，说法正确的是（    ）
%% answer: 
\jdanswer{
CD
}
%% memo:
\memoanswer{
A．“观察电容器的充、放电现象”实验中，充电时电流逐渐减小，放电时电流逐渐增大，故 A 错误；

B．“用油膜法估测油酸分子的大小”实验中，滴入油酸酒精溶液后，需待油酸溶液全部散开，形状稳定后，用一玻璃板轻轻盖在浅盘上，然后用水笔把油酸溶液的轮廓画在玻璃板上，测出油膜面积，故 B 错误；

C．“观察光敏电阻特性”实验中，光照强度增加，光敏电阻阻值减小；“观察金属热电阻特性”实验中，温度升高，金属热电阻阻值增大，故 C 正确；

D．“探究变压器原、副线圈电压与匝数的关系”实验中，如果可拆变压器的“横梁”铁芯没装上，原线圈接入 10 V 的交流电时，变压器的效果减弱，副线圈磁通量还是会发生变化，副线圈输出电压不为零，故 D 正确。

故选 CD。
}

\item
%% number: 17
%% paperName: 2023年6月浙江高考第 17 题 8 分
%% typeId: 6
%% type: 计算题
%% chapter: 气体性质
%% point: 热力学第一定律
%% method: 无
%% score: 8
%% degree: 650
%% duplicateId: 0
%% body:
如图所示，导热良好的固定直立圆筒内用面积 $S$ = 100 cm$^{2}$，质量 $m$ = 1 kg 的活塞封闭一定质量的理想气体，活塞能无摩擦滑动。圆筒与温度 300 K 的热源接触，平衡时圆筒内气体处于状态 A，其体积 $V_{A}$ = 600 cm$^{3}$。缓慢推动活塞使气体达到状态 B，此时体积 $V_{B}$ = 500 cm$^{3}$。固定活塞，升高热源温度，气体达到状态 C，此时压强 $p_{C}$ = 1.4$\times$10$^{5}$ Pa。已知从状态 A 到状态 C，气体从外界吸收热量 $Q$ = 14 J；从状态 B 到状态 C，气体内能增加 $\Delta U$ = 25 J；大气压 $p_{0}$ = 1.01$\times$10$^{5}$ Pa。
\begin{enumerate}
	\item
	气体从状态 A 到状态 B，其分子平均动能\_\_\_\_\_\_\_\_（选填“增大”、“减小”或“不变”），圆筒内壁单位面积受到的压力\_\_\_\_\_\_\_\_（选填“增大”、“减小”或“不变”）；
	\item
	求气体在状态 C 的温度 $T_{c}$；
	\item
	求气体从状态 A 到状态 B 过程中外界对系统做的功 $W$。
\end{enumerate}
\onepicture[align=right]{figs/17.png}
%% answer: 
\jdanswer{
\begin{enumerate}
	\item
	不变；增大

	\item
	350 K

	\item
	11 J

\end{enumerate}
}
%% memo:
\memoanswer{
（1）圆筒导热良好，则气体从状态 A 缓慢推动活塞到状态 B，气体温度不变，则气体分子平均动能不变；气体体积减小，则压强变大，圆筒内壁单位面积受到的压力增大；

（2）状态 A 时的压强 $p_{A} = p_{0} - \frac{{mg}}{S}$ = 1.0$\times$10$^{5}$ Pa

温度 $T_{A}$ = 300 K；体积 $V_{A}$ = 600 cm$^{3}$；

C 态压强 $p_{C}$ = 1.4$\times$10$^{5}$ Pa；温度 $T_{C}$ = ？；体积 $V_{C}$ = 500 cm$^{3}$；

根据 $\frac{{{p_{\mathrm{A}}}{V_{\mathrm{A}}}}}{{{T_{\mathrm{A}}}}} = \frac{{{p_{\mathrm{C}}}{V_{\mathrm{C}}}}}{{{T_{\mathrm{C}}}}}$

解得 $T_{C}$ = 350 K

（3）从 B 到 C 气体进行等容变化，则 $W_{BC}$ = 0，因从 B 到 C 气体内能增加 25 J 可知，气体从外界吸热 25 J，而气体从 A 到 C 从外界吸热 14 J，可知气体从 A 到 B 气体放热 11 J，从 A 到 B 气体内能不变，可知从 A 到 B 外界对气体做功 11 J。
}

\item
%% number: 18
%% paperName: 2023年6月浙江高考第 18 题 11 分
%% typeId: 6
%% type: 计算题
%% chapter: 运动和力
%% point: 动量和能量
%% method: 无
%% score: 11
%% degree: 450
%% duplicateId: 0
%% body:
为了探究物体间碰撞特性，设计了如图所示的实验装置。水平直轨道 AB、CD 和水平传送带平滑无缝连接，两半径均为 $R$ = 0.4 m 的四分之一圆周组成的竖直细圆弧管道 DEF 与轨道 CD 和足够长的水平直轨道 FG 平滑相切连接。质量为 3$m$ 的滑块 b 与质量为 2$m$ 的滑块 c 用劲度系数 $k$ = 100 $\UNm$ 的轻质弹簧连接，静置于轨道 FG 上。现有质量 $m$ = 0.12 kg 的滑块 a 以初速度 $v_{0}$ = 2$\sqrt {21} $ $\Ums$ 从 D 处进入，经 DEF 管道后，与 FG 上的滑块 b 碰撞（时间极短）。已知传送带长 $L$ = 0.8 m，以 $v$ = 2 $\Ums$ 的速率顺时针转动，滑块a与传送带间的动摩擦因数 $\mu$ = 0.5，其它摩擦和阻力均不计，各滑块均可视为质点，弹簧的弹性势能 $E_{p} = \frac{1}{2}kx^{2}$（$x$ 为形变量）。
\begin{enumerate}
	\item
	求滑块 a 到达圆弧管道 DEF 最低点 F 时速度大小 $v_{F}$ 和所受支持力大小 $F_{N}$；
	\item
	若滑块 a 碰后返回到 B 点时速度 $v_{B}$ = 1 $\Ums$，求滑块 a、b 碰撞过程中损失的机械能$\Delta E$；
	\item
	若滑块 a 碰到滑块 b 立即被粘住，求碰撞后弹簧最大长度与最小长度之差 $\Delta x$。
\end{enumerate}
\onepicture[align=right]{figs/18.png}
%% answer: 
\jdanswer{
\begin{enumerate}
	\item
	10 $\Ums$；31.2 N

	\item
	0

	\item
	0.2 m

\end{enumerate}
}
%% memo:
\memoanswer{
（1）滑块 a 从 D 到 F，由能量关系

 $mg \cdot 2R = \frac{1}{2}mv_{F}^{2} - \frac{1}{2}mv_{0}^{2}$

在 F 点  $F_{N}-mg=m\frac{v_{F}^{2}}{R}$

解得  $v_{F}=10\ \Ums,\ F_{N}=31.2\ N$

（2）滑块 a 返回 B 点时的速度  $v_{B}=1\ \Ums$，滑块 a 一直在传送带上减速，加速度大小为  $a=\mu g=5\ \Ums^{2}$

根据  $v_{B}^{2}=v_{C}^{2}-2aL$ 可得在 C 点的速度  $v_{C}=3\ \Ums$。

则滑块 a 从碰撞后到到达 C 点

 $\frac{1}{2}mv_{1}^{2}=\frac{1}{2}mv_{C}^{2}+mg\cdot2R$

解得  $v_{1}=5\ \Ums$

因 ab 碰撞动量守恒，则

 $mv_{F}=-mv_{1}+3mv_{2}$

解得碰后 b 的速度  $v_{2}=5\ \Ums$

则碰撞损失的能量

$$ \Delta E=\frac{1}{2}mv_{F}^{2}-\frac{1}{2}mv_{1}^{2}-\frac{1}{2}\cdot3mv_{2}^{2}=0 $$

（3）若滑块 a 碰到滑块 b 立即被粘住，则 ab 碰后的共同速度  $m v_{F} = 4m v$

解得  $v = 2.5 \, \Ums$

当弹簧被压缩到最短或者伸长到最长时有共同速度

 $4mv = 6mv'$

则  $v' = \frac{5}{3} \Ums$

当弹簧被压缩到最短时压缩量为  $x_{1}$，由能量关系

 $\frac{1}{2} \cdot 4mv^{2} = \frac{1}{2} \cdot 6mv'^{2} + \frac{1}{2} kx_{1}^{2}$

解得  $x_{1} = 0.1 \, m$

同理当弹簧被拉到最长时伸长量为  $x_{2} = x_{1}$

则弹簧最大长度与最小长度之差  $\Delta x = 2x_{1} = 0.2 \, m$
}

\item
%% number: 19
%% paperName: 2023年6月浙江高考第 19 题 11 分
%% typeId: 6
%% type: 计算题
%% chapter: 电磁感应
%% point: 电磁感应能量分析
%% method: 无
%% score: 11
%% degree: 450
%% duplicateId: 0
%% body:
某兴趣小组设计了一种火箭落停装置，简化原理如图所示，它由两根竖直导轨、承载火箭装置（简化为与火箭绝缘的导电杆 MN）和装置 A 组成，并形成团合回路。装置 A 能自动调节其输出电压确保回路电流 $I$ 恒定，方向如图所示。导轨长度远大于导轨间距，不论导电杆运动到什么位置，电流 $I$ 在导电杆以上空间产生的磁场近似为零，在导电杆所在处产生的磁场近似为匀强磁场，大小 $B_{1} = kI$（其中 $k$ 为常量），方向垂直导轨平面向里；在导电杆以下的两导轨间产生的磁场近似为匀强磁场，大小 $B_{2}$ = 2$kI$，方向与 $B_{1}$ 相同。火箭无动力下降到导轨顶端时与导电杆粘接，以速度 $v_{0}$ 进入导轨，到达绝缘停靠平台时速度恰好为零，完成火箭落停。已知火箭与导电杆的总质量为 $M$，导轨间距 $d = \frac{{3Mg}}{{k{I^2}}}$，导电杆电阻为 $R$。导电杆与导轨保持良好接触滑行，不计空气阻力和摩擦力，不计导轨电阻和装置 A 的内阻。在火箭落停过程中，
\begin{enumerate}
	\item
	求导电杆所受安培力的大小 $F$ 和运动的距离 $L$；
	\item
	求回路感应电动势 $E$ 与运动时间 $t$ 的关系；
	\item
	求装置 A 输出电压 $U$ 与运动时间 $t$ 的关系和输出的能量 $W$；
	\item
	若 $R$ 的阻值视为 0，装置 A 用于回收能量，给出装置 A 可回收能量的来源和大小。
\end{enumerate}
\onepicture[align=right]{figs/19.png}
%% answer: 
\jdanswer{
\begin{enumerate}
	\item
	3$Mg$；$\frac{{v_0^2}}{{4g}}$

	\item
	$E = \frac{{6Mg}}{I}$($v_{0} -$ 2$gt$)

	\item
	$U = \frac{{3Mg}}{I}$($v_{0} -$ 2$gt$) + $IR$；$W$ = 3$Mg$($v_{0}t -$ 2$gt^{2}$) + $I^{2}Rt$

	\item
	装置 A 可回收火箭的动能和重力势能；$\frac{3}{4}Mv_{0}^{2}$

\end{enumerate}
}
%% memo:
\memoanswer{
（1）导体杆受安培力 $F = B_{1}Id$ = 3$Mg$

方向向上，则导体杆向下运动的加速度 $Mg - F = Ma$

解得 $a = -$ 2$g$

导体杆运动的距离 $L = \frac{{0 - v_0^2}}{{2a}} = \frac{{v_0^2}}{{4g}}$

（2）回路的电动势 $E = B_{2}dv$

其中 $v = v_{0} + at$

解得 $E = \frac{{6Mg}}{I}$($v_{0} -$ 2$gt$)

（3）由能量关系 $UI = Fv + I^{2}R$

其中 $v = v_{0} + at$

可得 $U = \frac{{3Mg}}{I}$ ($v_{0} -$ 2$gt$) + $IR$

输出能量 $W = UIt$ = 3$Mg$($v_{0}t -$ 2$gt^{2}$) + $I^{2}Rt$

（4）装置 A 可回收火箭的动能和重力势能；从开始火箭从速度 $v_{0}$ 到平台速度减为零，则

${{E' = \frac{1}{2}Mv_{0}}^{2} + MgL = \frac{3}{4}Mv_{0}}^{2}$
}

\item
%% number: 20
%% paperName: 2023年6月浙江高考第 20 题 11 分
%% typeId: 6
%% type: 计算题
%% chapter: 磁场
%% point: 带电粒子在磁场中的运动
%% method: 无
%% score: 11
%% degree: 350
%% duplicateId: 0
%% body:
利用磁场实现离子偏转是科学仪器中广泛应用的技术。如图所示，$Oxy$ 平面（纸面）的第一象限内有足够长且宽度均为 $L$、边界均平行 $x$ 轴的区域 Ⅰ 和 Ⅱ，其中区域存在磁感应强度大小为 $B_{1}$ 的匀强磁场，区域 Ⅱ 存在磁感应强度大小为 $B_{2}$ 的磁场，方向均垂直纸面向里，区域 Ⅱ 的下边界与 $x$ 轴重合。位于（0，3$L$）处的离子源能释放出质量为 $m$、电荷量为 $q$、速度方向与 $x$ 轴夹角为 60${^\circ}$ 的正离子束，沿纸面射向磁场区域。不计离子的重力及离子间的相互作用，并忽略磁场的边界效应。
\begin{enumerate}
	\item
	求离子不进入区域Ⅱ的最大速度 $v_{1}$ 及其在磁场中的运动时间$t$；
	\item
	若 $B_{2}$ = 2$B_{1}$，求能到达 $y = \frac{L}{2}$ 处的离子的最小速度 $v_{2}$；
	\item
	若 $B_{2} = \frac{{{B_1}}}{L}y$，且离子源射出的离子数按速度大小均匀地分布在 $\frac{{{B_1}qL}}{m}$ \textasciitilde{} $\frac{{6{B_1}qL}}{m}$ 范围，求进入第四象限的离子数与总离子数之比 $\eta$。
\end{enumerate}
\onepicture[align=right]{\includesvg{figs/20}}
%% answer: 
\jdanswer{
\begin{enumerate}
	\item
	$v_{1} = \frac{{2{B_1}qL}}{m}$；$t = \frac{{\pi m}}{{3q{B_1}}}$

	\item
	$v_{2} = \frac{{4{B_1}qL}}{m}$

	\item
	60\%

\end{enumerate}
}
%% memo:
\memoanswer{
（1）当离子不进入磁场Ⅱ速度最大时，轨迹与边界相切，则由几何关系

 $r_{1}\cos60^{\circ}=r_{1}-L$

解得  $r_{1}=2L$

根据  $qv_{1}B_{1}=m\frac{v_{1}^{2}}{r_{1}}$，解得  $v_{1}=\frac{2B_{1}qL}{m}$

在磁场中运动的周期  $T = \frac{2\pi m}{qB_{1}}$，运动时间  $t =$

 $\frac{60^{\circ}}{360^{\circ}}T=\frac{\pi m}{3qB_{1}}$

（2）若  $B_{2}=2B_{1}$，根据

 $r=\frac{mv}{qB}$

可知  $r_{1}=2r_{2}$

粒子在磁场中运动轨迹如图，设  $O_{1}O_{2}$ 与磁场边界夹角为  $\alpha$，由几何关系

 $r_{1}\sin\alpha-r_{1}\sin30^{\circ}=L$

 $r_{2}-r_{2}\sin\alpha=\frac{L}{2}$

解得  $r_{2}=2L$， $\sin\alpha=\frac{3}{4}$

根据  $qv_{2}B_{2}=m\frac{v_{2}^{2}}{r_{2}}$

解得  $v_{2}=\frac{4B_{1}qL}{m}$

（3）当最终进入区域Ⅱ的粒子若刚好到达 x 轴，则由动量定理  $B_{2}qv_{y}\Delta t = m\Delta v_{x}$，即  $\frac{B_{1}}{L}yq\Delta y = m\Delta v_{x}$，求和可得  $\sum\frac{B_{1}}{L}yq\Delta y = \sum m\Delta v_{x}$

粒子从区域Ⅰ到区域Ⅱ最终到 x 轴上的过程中

 $m(v-v\cos60^{\circ})=B_{1}qL+\frac{B_{1}}{L}\cdot\frac{0+L}{2}\cdot qL$

解得  $v=\frac{3B_{1}qL}{m}$

则速度在  $\frac{3B_{1}qL}{m} \sim \frac{6B_{1}qL}{m}$ 之间的粒子才能进入第四象限；因离子源射出粒子的速度范围在  $B_{1}qL \sim 6B_{1}qL$，又粒子源射出的粒子个数按速度大小均匀分布，可知能进入第四象限的粒子占粒子总数的比例为  $\eta = 60\%$
}

\end{enumerate}

\end{document}
